\documentclass{beamer}
\usepackage[utf8]{inputenc}

\usetheme{Madrid}
\usecolortheme{default}
\usepackage{amsmath,amssymb,amsfonts,amsthm}
\usepackage{txfonts}
\usepackage{tkz-euclide}
\usepackage{listings}
\usepackage{adjustbox}
\usepackage{array}
\usepackage{tabularx}
\usepackage{gvv}
\usepackage{lmodern}
\usepackage{circuitikz}
\usepackage{tikz}
\usepackage{graphicx}
\usepackage[T1]{fontenc}
\usepackage[utf8]{inputenc}

% Settings for C and Python code listings
\lstset{
  language=Python,
  basicstyle=\ttfamily\small,
  breaklines=true,
  showstringspaces=false,
  numbers=left,
  numberstyle=\tiny\color{gray},
  keywordstyle=\color{blue},
  stringstyle=\color{orange},
  commentstyle=\color{green!60!black},
  literate={->}{{$\rightarrow$}}1
}

\setbeamertemplate{page number in head/foot}[totalframenumber]

\begin{document}

\title{2.4.6}
\date{September 22, 2025}
\author{Taraka Abhinav - EE25BTECH11016}

\frame{\titlepage}

% --- Mathematical Solution ---

\begin{frame}{Question}
If $\vec{a} = 5\hat{i} - \hat{j} - 3\hat{k}$ and $\vec{b} = \hat{i} + 3\hat{j} - 5\hat{k}$, show that the vectors $\vec{a} + \vec{b}$ and $\vec{a} - \vec{b}$ are perpendicular.
\end{frame}

\begin{frame}{Solution - Part 1}
Two vectors are perpendicular if their dot product is zero. We need to show that $(\mathbf{a} + \mathbf{b})^\top (\mathbf{a} - \mathbf{b}) = 0$.

First, we find the vectors $\vec{a} + \vec{b}$ and $\vec{a} - \vec{b}$.
The given vectors can be written in column form as:
\begin{equation}
\mathbf{a} = \begin{pmatrix} 5 \\ -1 \\ -3 \end{pmatrix}, \quad \mathbf{b} = \begin{pmatrix} 1 \\ 3 \\ -5 \end{pmatrix}
\end{equation}
\end{frame}

\begin{frame}{Solution - Part 2}
Now, we calculate the sum and difference:
\begin{equation}
\mathbf{a} + \mathbf{b} = \begin{pmatrix} 5+1 \\ -1+3 \\ -3-5 \end{pmatrix} = \begin{pmatrix} 6 \\ 2 \\ -8 \end{pmatrix}
\end{equation}
\begin{equation}
\mathbf{a} - \mathbf{b} = \begin{pmatrix} 5-1 \\ -1-3 \\ -3-(-5) \end{pmatrix} = \begin{pmatrix} 4 \\ -4 \\ 2 \end{pmatrix}
\end{equation}
Next, we compute the dot product of these two resultant vectors:
\begin{align}
(\mathbf{a} + \mathbf{b})^\top (\mathbf{a} - \mathbf{b}) &= \begin{pmatrix} 6 & 2 & -8 \end{pmatrix} \begin{pmatrix} 4 \\ -4 \\ 2 \end{pmatrix} \nonumber \\
&= (6)(4) + (2)(-4) + (-8)(2) \nonumber \\
&= 24 - 8 - 16 \nonumber \\
&= 0
\end{align}
\end{frame}

\begin{frame}{Conclusion}
Since the dot product is zero, the vectors $(\vec{a} + \vec{b})$ and $(\vec{a} - \vec{b})$ are perpendicular. Hence, proved.
\end{frame}

% --- Code Implementation ---

\begin{frame}[fragile]
    \frametitle{C Code - Calculation Function}
    \lstset{language=C}
    \begin{lstlisting}
// File: vector_ops.c
float check_perpendicularity(
    float ax, float ay, float az, 
    float bx, float by, float bz) {
    
    // Calculate p = a + b
    float px = ax + bx;
    float py = ay + by;
    float pz = az + bz;

    // Calculate q = a - b
    float qx = ax - bx;
    float qy = ay - by;
    float qz = az - bz;

    // Return the dot product of p and q
    return (px * qx) + (py * qy) + (pz * qz);
}
    \end{lstlisting}
\end{frame}

\begin{frame}[fragile]
    \frametitle{Python Code - Setup and C Library Import}
    \lstset{language=Python}
    \begin{lstlisting}
import numpy as np
import matplotlib.pyplot as plt
import ctypes
import os

# Load the compiled C library
lib_name = 'vector_ops.so'
c_lib = ctypes.CDLL(os.path.abspath(lib_name))

# Define C function argument and return types
c_lib.check_perpendicularity.argtypes = [ctypes.c_float] * 6
c_lib.check_perpendicularity.restype = ctypes.c_float
    \end{lstlisting}
\end{frame}

\begin{frame}[fragile]
    \frametitle{Python Code - Vector Definition and Calculation}
    \lstset{language=Python}
    \begin{lstlisting}
# Define vectors and call C function
a = np.array([5, -1, -3], dtype=np.float32)
b = np.array([1, 3, -5], dtype=np.float32)
p = a + b
q = a - b

# Call the C function to get the dot product
dot_product = c_lib.check_perpendicularity(*a, *b)
print(f"Dot product from C: {dot_product}")
if abs(dot_product) < 1e-6:
    print("Vectors are perpendicular.")
    \end{lstlisting}
\end{frame}

\begin{frame}[fragile]
    \frametitle{Python Code - Plotting Logic}
    \lstset{language=Python}
    \begin{lstlisting}
# Generate the 3D plot
fig = plt.figure(figsize=(10, 8))
ax = fig.add_subplot(111, projection='3d')
# Dictionary of vectors to plot
vectors = {
    '$\\vec{a}$': (a, 'green'),
    '$\\vec{b}$': (b, 'orange'),
    '$\\vec{a}+\\vec{b}$': (p, 'blue'),
    '$\\vec{a}-\\vec{b}$': (q, 'red')
}
origin = [0, 0, 0]

for name, (vec, color) in vectors.items():
    ax.quiver(*origin, *vec, color=color, 
              arrow_length_ratio=0.1, label=name)
    ax.scatter(*vec, color=color, s=50)
    ax.text(vec[0], vec[1], vec[2], 
            f'({vec[0]},{vec[1]},{vec[2]})')
    \end{lstlisting}
\end{frame}

\begin{frame}[fragile]
    \frametitle{Python Code - Final Plotting Details}
    \lstset{language=Python}
    \begin{lstlisting}
# Plotting settings
ax.set_xlabel('X Axis')
ax.set_ylabel('Y Axis')
ax.set_zlabel('Z Axis')
ax.set_title('3D Plot of Vectors and Endpoints')
ax.legend()
ax.grid(True)
ax.view_init(elev=20, azim=25) # Set view angle
plt.show()
    \end{lstlisting}
\end{frame}


\begin{frame}{Generated 3D Plot}
\begin{figure}
    \centering
    \includegraphics[width=1\linewidth]{figs/Figure_1.png}
    \label{fig:placeholder}
\end{figure}
\end{frame}

\end{document}